\chapter*{How to Read This Book}
\addcontentsline{toc}{chapter}{How to Read This Book}
\markboth{HOW TO READ THIS BOOK}{HOW TO READ THIS BOOK}

Each chapter starts from a real situation, the kind you find on a whiteboard
when a system is about to grow, and asks why the technologies we reach for are
shaped the way they are.

The answer is always historical. No idea in computing showed up because it was
elegant. Each one was forced by a pain somebody had, at a certain time and
place. It fixed something its predecessor got wrong, and it brought new costs of
its own. If you know that lineage, you know when to use the idea and when to run
from it. If you only know the rule, you'll apply it where the pain doesn't exist.

\section*{Two kinds of boxes}

Stories and concepts are kept apart, so you always know which one you're
reading.

\begin{history}{year \textperiodcentered{} place}
An amber box tells a story: who did it, when, and what forced the idea into
existence. Everything in it is real: the people, the dates and the problems.
No invented incidents.
\end{history}

\begin{concept}{name}
A blue box defines an idea you'll need again. It always ends the same way:

\textbf{What it fixes:} the pain that made the idea worth having.

\textbf{What it costs:} what you give up by using it.
\end{concept}

The text between the boxes is the argument that connects them.

\section*{Diagrams}

Colors in diagrams are roles, not decoration. A color means the same thing in
every figure of the book.

\begin{center}
\begin{tikzpicture}
  \matrix[column sep=6mm, row sep=4mm, nodes={anchor=center}] {
    \node[hf/source, minimum width=26mm] {source}; &
    \node[hf/note, anchor=west, text width=95mm, align=left] {where the data comes from: a plant, a device, a sensor}; \\
    \node[hf/reader, minimum width=26mm] {reader}; &
    \node[hf/note, anchor=west, text width=95mm, align=left] {a program that consumes data: a screen, a client, an app}; \\
    \node[hf/store, minimum width=26mm] {store}; &
    \node[hf/note, anchor=west, text width=95mm, align=left] {where data is kept: a historian, a database, a copy}; \\
    \node[hf/bus, minimum width=26mm] {middle}; &
    \node[hf/note, anchor=west, text width=95mm, align=left] {something that sits between: a broker, a standard, a bus}; \\
    \draw[hf/ask] (-13mm,0) -- (13mm,0); &
    \node[hf/note, anchor=west, text width=95mm, align=left] {asks: a request, a poll}; \\
    \draw[hf/answer] (-13mm,0) -- (13mm,0); &
    \node[hf/note, anchor=west, text width=95mm, align=left] {answers a request}; \\
    \draw[hf/push] (-13mm,0) -- (13mm,0); &
    \node[hf/note, anchor=west, text width=95mm, align=left] {pushes: sent without being asked}; \\
    \draw[hf/save] (-13mm,0) -- (13mm,0); &
    \node[hf/note, anchor=west, text width=95mm, align=left] {keeps: written to storage}; \\
    \draw[hf/pain, dashed] (-13mm,0) -- (13mm,0); &
    \node[hf/note, anchor=west, text width=95mm, align=left, text=hfAlert] {a cost, a pain, or something forbidden}; \\
  };
\end{tikzpicture}
\end{center}

\section*{Exact facts, firm opinions}

This isn't a scientific book, and it doesn't pretend to be one.

The history is checked. When the book says who did something and when, that's
what happened.

The opinions are another matter. When something is a mistake, the book says
it's a mistake. When something is useless, it says \emph{useless}, not ``of
limited value in some contexts''. Sometimes that will be unfair or imprecise, and
we accept that.

A reader who leaves with a position can argue with it, refine it, or throw it
away. A reader who leaves with ``it depends'' has nothing to hold. Between being
precise and leaving you with a position, this book prefers to give you a position.
